\clearpage
\section{Retained residue-counting continuation}
\label{sec:historical-residue-continuation}
The following arguments develop residue counting and the associated norm-image conditions. Their formulas and hypotheses remain distinct from the amended conditional completion in Section~\ref{sec:conditional-final}. The formal companion distinguishes the source contracts and their verification status.
\providecommand{\OO}{\mathcal O}
\providecommand{\ZZ}{\mathbb Z}
\providecommand{\RR}{\mathbb R}
\providecommand{\CC}{\mathbb C}
\providecommand{\NN}{\mathbb N}
\providecommand{\Cl}{\operatorname{Cl}}
\providecommand{\Nm}{\operatorname N}
\providecommand{\vol}{\operatorname{vol}}
\providecommand{\covol}{\operatorname{covol}}
\providecommand{\absn}[1]{\lvert #1\rvert}




\paragraph{Scope and notation.}
Let $K$ be a number field of degree $d\geq1$, with ring of integers
$\OO=\OO_K$. Write $\Nm I=\#(\OO/I)$ for the absolute norm of a nonzero
integral ideal and $n(x)=\Nm_{K/\mathbb Q}(x)\in\ZZ$ for $x\in\OO$.
Thus $\Nm(x)=\absn{n(x)}$. Let $\iota$ be the mixed embedding in its
Euclidean coordinates, $L(J)=\iota(J)$, and $D$ the actual unit-norm
fundamental cone used in the preceding counting block. Its norm function
scales by $t^d$ under a positive real dilation $t$.
The full-paper ledger remains \textbf{17 checked claims and 3 partial claims}.

\paragraph{Actual additive residues.}
Fix nonzero integral ideals $M\subseteq J$, and put $q=[J:M]=\#(J/M)$.
The formal quotient of ideal lattices is identified with the additive
congruence relation $x-a\in M$ on actual elements of $J$. In particular,
the residue label is not an independent parameter attached to a count.

\begin{proposition}\label{prop:residue-count}
There is a constant $B\geq0$, depending on $K,M,J$, such that, for every
$a\in J/M$ and every real $t\geq1$,
\[
 \left|\#\{x\in J:x\bmod M=a,\ \iota(x)\in tD\}
 -\frac{\vol(D)}{q\,\covol(L(J))}t^d\right|
 \leq B(t^{d-1}+1).
\]
The same constant works for all residues $a$.
\end{proposition}

The estimate is proved from the actual boundary charts and the fixed-lattice
discrepancy theorem. The index factor comes from the covolume identity
$\covol(L(M))=q\,\covol(L(J))$.
For a complex weight $w:J/M\to\CC$, summing the estimates gives the leading
term
\[
 \frac{\vol(D)t^d}{q\,\covol(L(J))}\sum_{a\in J/M}w(a),
\]
with error at most $B\sum_a\absn{w(a)}(t^{d-1}+1)$.
If the finite sum of weights is zero, only that error remains.
Writing $Y=t^d$ gives an error $O(Y^{1-1/d}+1)$ for the actual norm threshold.
These assertions include $d=1$.

\smallskip
\noindent\emph{Lean attachment:}
\nolinkurl{ResidueLatticeIdeal}, \nolinkurl{ResidueLatticeIdealPoints},
\nolinkurl{ResidueLatticeIdealWeights}, and \nolinkurl{ResidueLatticeNormCounts}.

\newpage
\paragraph{The residue norm and its sign.}
For $m\in\NN$, reducing the multiplication matrices of $\OO$ modulo $m$
and taking their determinants constructs a monoid homomorphism
\[
 \overline n:\OO/m\OO\longrightarrow\ZZ/m\ZZ,
 \qquad \overline n([x])=n(x)\bmod m.
\]
The construction is proved independent of the integral basis.
For $m\ne0$, its source is finite of cardinality $m^d$.
Consequently $x-y\in m\OO$ implies $n(x)\equiv n(y)\pmod m$.
No surjectivity of $\overline n$ is asserted.

The relation to the geometric sign is also proved. If $\sigma$ ranges over
real embeddings and $\tau$ over one embedding from each complex pair, then
\[
 n(x)=\prod_{\sigma}\sigma(x)\prod_{\tau}\absn{\tau(x)}^2.
\]
For $x\ne0$, the second factor is positive. Thus the sign of $n(x)$ is
exactly the sign of the real-coordinate product. Cutting $D$ into its
positive and negative portions $D_+,D_-$ adds no boundary outside the
existing frontier, and both portions retain the proved boundary atlas.
They are disjoint and their union is $D$.

For $\varepsilon\in\{1,-1\}$ define
$\overline n_{\varepsilon}(a)=\varepsilon\overline n(a)$.
On the corresponding sign portion,
\[
 \Nm(x)\bmod m=\overline n_{\varepsilon}([x]).
\]
The sign cannot be dropped: $-1\equiv2\pmod3$, whereas
$\Nm(-1)=1$ and $\Nm(2)=2$ for principal ideals of $\ZZ$.

\paragraph{The exact ideal filter.}
Suppose $IJ=(x)$, $n(x)$ has the indicated sign, and
$\gcd(\Nm J,m)=1$. For $c\in\ZZ/m\ZZ$,
\[
 \Nm I\equiv c\pmod m
 \quad\Longleftrightarrow\quad
 \overline n_{\varepsilon}([x])=c(\Nm J\bmod m).
\]
This gives an actual finite residue target, which is allowed to be empty.
The product identity before cancellation does not require coprimality.

Let $[J]=C^{-1}$ in $\Cl(\OO)$, and let $\mathcal G(J)$ denote the
actual fundamental-cone generator set. For every predicate $P:\RR\to\mathrm{Prop}$
and every real $X$, the norm-preserving correspondence gives
\begin{align*}
 &\#\{x\in\mathcal G(J):\Nm(x)\leq X\Nm J,\ P(\Nm(x))\}\\
 &\qquad=\#\mu(K)\,
 \#\{I\ne0:[I]=C,\ \Nm I\leq X,\ P(\Nm I\,\Nm J)\}.
\end{align*}
The same norm-preserving equivalence carries arbitrary complex weights.
In particular, write $G_{J,\chi}(Y)$ for the sum of
$\chi(\Nm(x)\bmod m)$ over these generators with $\Nm(x)\leq Y$, and
$T_{C,\chi}(X)$ for the corresponding sum over ideals in class $C$.
For multiplicative $\chi$ preserving $0$ and $1$, the exact identity is
\begin{equation}\label{eq:weighted-torsion}
 G_{J,\chi}(X\Nm J)
 =\#\mu(K)\,\chi(\Nm J\bmod m)\,T_{C,\chi}(X).
\end{equation}
Both signs are included before applying this full roots-of-unity factor.
A root of unity can interchange the sign portions, so that factor is not
automatically valid for either portion separately.

\smallskip
\noindent\emph{Lean attachment:}
\nolinkurl{NormResidueArithmetic}, \nolinkurl{NormResidueIdealFilter},
\nolinkurl{NormSignArithmetic}, \nolinkurl{IdealNormFilterAttachment}, and
\nolinkurl{IdealNormWeightedAttachment}.

\newpage
\paragraph{Ordinary class characters: an unconditional power saving.}
The preceding block proved, for every ordinary ideal class $C$,
\[
 A_C(X)=c_KX+O_{K,C}(X^{1-1/d}),\qquad
 c_K=\frac{2^{r_1}(2\pi)^{r_2}R_K}{\#\mu(K)\sqrt{\absn{\Delta_K}}}.
\]
Here $A_C$ counts actual nonzero integral ideals, $r_1,r_2$ count the real
and complex places, $R_K$ is the regulator, and $\Delta_K$ is the discriminant. For a weight
$w:\Cl(\OO)\to\CC$, define the actual finite ideal sum
\[
 S_w(X)=\sum_{\substack{I\ne0\\\Nm I\leq X}}w([I]).
\]
Exact grouping by ideal class gives $S_w(X)=\sum_C w(C)A_C(X)$.
If $\sum_Cw(C)=0$, the common leading coefficient cancels, so there is
$B_w\geq0$ such that
\[
 \absn{S_w(X)}\leq B_wX^{1-1/d}\qquad(X\geq1).
\]
For every nontrivial character $\chi:\Cl(\OO)\to\CC^{\times}$,
the finite character sum is zero. Thus the displayed estimate applies
without an additional arithmetic cancellation hypothesis.
This conclusion concerns ordinary class characters. It does not identify
them with the further norm-residue characters needed for Frobenius separation.

\paragraph{An actual Mellin endpoint.}
The counting functions are constant on the integer norm steps. Their
measurability, and hence that of $S_w$, is proved from this fact.
For zero-sum $w$, define
\[
 F_w(h)=(1+h)\int_1^{\infty}x^{-2-h}S_w(x)\,dx
 \qquad(h\geq0).
\]
Since $1-1/d<1$, the logarithmic first moment
\[
 M_w=\int_1^{\infty}(1+\log x)\,\absn{x^{-2}S_w(x)}\,dx
\]
is finite. Lean proves the explicit estimate
$\absn{F_w(h)-F_w(0)}\leq hM_w$ for every $h\geq0$, and therefore
$F_w(h)=F_w(0)+O(h)$ as $h\downarrow0$.
The function here is the displayed integral of the actual ideal sum.
Its identification with the intended arithmetic Euler product is a separate step.

For a finite family of nontrivial ordinary class characters, suppose functions
$p,Z$ satisfy $h\,p(h)\to P\ne0$, $h\,Z(h)\to Q\ne0$, and, for all sufficiently
small $h>0$,
\[
 Z(h)=\left(p(h)\prod_i F_{\chi_i}(h)\right)^e,
 \qquad e\in\NN.
\]
The checked real-axis criterion then gives $e=1$ and
$F_{\chi_i}(0)\ne0$ for every $i$. The pole limits and product identity
remain hypotheses of this application.

\smallskip
\noindent\emph{Lean attachment:}
\nolinkurl{IdealClassCharacterBound}, \nolinkurl{IdealCharacterMellinMeasurable},
\nolinkurl{IdealCharacterMellin}, and \nolinkurl{IdealCharacterMellinPole}.

\newpage
\paragraph{An actual unit action on ideal residues.}
Fix $m>0$ and a nonzero integral ideal $J$.
Multiplication by $u\in(\OO/m\OO)^{\times}$ gives a proved permutation of
$J/mJ$. Its inverse follows from $(ab-1)x\in mJ$ when
$ab-1\in m\OO$ and $x\in J$. No surjectivity of
$J/mJ\to\OO/m\OO$ or coprimality condition on $J$ is used.
The well-defined residue norm on this quotient satisfies
\[
 \overline n_J(u\cdot a)=\overline n(u)\,\overline n_J(a).
\]
For a multiplicative $\chi:\ZZ/m\ZZ\to\CC$ preserving $1$, an exhibited
unit with
\begin{equation}\label{eq:unit-witness}
 u\in(\OO/m\OO)^{\times},\qquad \chi(\overline n(u))\ne1
\end{equation}
therefore forces
\[
 \sum_{a\in J/mJ}\chi\bigl(s\,\overline n_J(a)\bigr)=0
 \qquad(s\in\ZZ/m\ZZ).
\]
Taking $s=1,-1$ gives the two finite means needed by the sign-corrected
norm weights. Their vanishing follows from the actual permutation.

\paragraph{The assembled ideal norm-character estimate.}
\begin{proposition}\label{prop:norm-character}
Let $C\in\Cl(\OO)$, let $J\ne0$ be an integral ideal with $[J]=C^{-1}$,
and let $m>0$. Let $\chi:\ZZ/m\ZZ\to\CC$ preserve multiplication,
$0$, and $1$. Assume a residue-ring unit satisfying
\eqref{eq:unit-witness}, and assume
$\chi(\Nm J\bmod m)\ne0$. Then
\[
 T_{C,\chi}(X)
 =\sum_{\substack{I\ne0,\ [I]=C\\\Nm I\leq X}}
       \chi(\Nm I\bmod m)
 =O_{K,J,m,\chi}(X^{1-1/d}).
\]
More precisely, some $B\geq0$ bounds its absolute value by
$BX^{1-1/d}$ for every real $X\geq1$.
\end{proposition}

The assembly identifies the residue sum with the sum over actual generators,
recombines the sign portions, and uses $Y=t^d$ to obtain
$\absn{G_{J,\chi}(Y)}\leq B_0(Y^{1-1/d}+1)$.
Equation~\eqref{eq:weighted-torsion} then permits division by the nonzero
factor $\#\mu(K)\chi(\Nm J\bmod m)$.
The final theorem supplies its finite means and counting estimates internally.
For $d=1$ the exponent is zero and the conclusion is a bounded sum, so the
argument does not require $d>1$.

\paragraph{What remains open.}
Nontriviality of $\chi$ on the larger residue ring does not supply
\eqref{eq:unit-witness}: its restriction to the image of the norm on units
may be trivial. Witness existence for the intended characters, their
remaining norm-image cases, and suitable inverse-class representatives with
$\chi(\Nm J)\ne0$ are still arithmetic obligations.
Measurability and the Mellin attachment for these norm-character sums remain
to be checked. The compiled Mellin results above concern ordinary class characters.
The resulting estimates must also be attached to the required Euler-product
identities and prime-occurrence argument. The conditional pole criterion
does not supply those identities. Claims 9.1, 9.2, and 11.1 remain partial,
including the separate local compatibility obligations. The ledger stays
\textbf{17 checked / 3 partial}.

\smallskip
\noindent\emph{Lean attachment:}
\nolinkurl{IdealResidueUnitActionMean}, \nolinkurl{NormResidueUnitCounting},
\nolinkurl{NormResidueGeneratorAttachment}, \nolinkurl{IdealNormCharacterBound},
and \nolinkurl{NormResidueCharacterCounting}.
